\section{零点定理}\label{sec:nullstellensatz}

\BookPairedNote
  {本节叙述一个构成代数几何许多内容之基础的定理: 零点定理. 这个结果让我们能在仿射簇和多项式环中的根理想之间建立一种有力的对应.}
  {\textquotedbl{}Nullstellensatz\textquotedbl{} 是一个德语词, 由 \textquotedbl{}Nullstellen\textquotedbl{} (零点) 和 \textquotedbl{}Satz\textquotedbl{} (定理) 组成.}
\noindent 这种对应是连接几何与代数两界的词典的主干.

为了叙述这一定理, 需要对基域作一个关键假设.

\begin{Definition}{代数闭域}\leavevmode\label{def:algebraically-closed-field}
  如果 $K[x]$ 中每个非常数多项式在 $K$ 中都至少有一个零点, 则称域 $K$ 是\textbf{代数闭的}.
\end{Definition}

例如, 域 $\mathbb{Q}$ 和 $\mathbb{R}$ 都不是代数闭域, 因为 $x^{2}+1$ 没有有理零点, 也没有实零点. 然而, $x^{2}+1$ 在 $\mathbb{C}$ 中确实有零点, 即 $\mathrm{i}$ 和 $-\mathrm{i}$. 事实上, $\mathbb{C}[x]$ 中每个非常数多项式都至少有一个零点, 这就是代数基本定理的内容.

\begin{Theorem}{代数基本定理}\leavevmode\label{thm:fundamental-theorem-of-algebra}
  域 $\mathbb{C}$ 是代数闭的.
\end{Theorem}

虽然代数基本定理往往很早就讲到, 但它绝非显而易见. 它有很多证明, 包括利用 Galois 理论, 复分析和拓扑学的证明. 这些都不在本书的范围之内, 而 $\mathbb{C}$ 是代数闭域这一事实也不是本书内容在逻辑上所必需的. 这里引入代数基本定理, 只是为了强调至少存在一个熟悉而具体的代数闭域的例子. 我们鼓励读者不加证明地接受这一定理.

如我们在第~\ref{sec:vanishing-set-operator}~节所见, 基域是无限域时, 代数几何的许多内容更为直接, 因为这时不再需要区分多项式与其定义的函数. 在代数闭域上工作的一个好处, 是代数闭域必定是无限域.

\begin{Proposition}{代数闭域是无限域}\leavevmode\label{prop:algebraically-closed-fields-are-infinite}
  如果 $K$ 是代数闭域, 则 $K$ 是无限域.
\end{Proposition}

\begin{Proof}
  见 Exercise~\ref{ex:algebraically-closed-field-infinite}.
\end{Proof}

特别地, 对任何素数 $p\in\mathbb{Z}$, 有限域 $\mathbb{F}_{p}=\mathbb{Z}/\langle p\rangle$ 都不是代数闭域. 那么, 除了 $\mathbb{C}$ 以外, 哪些域是代数闭的呢? 也许再没有别的熟悉例子, 但有一种方法可以从任意域 $K$ 出发, 构造其\textbf{代数闭包} $\overline{K}$, 即包含 $K$ 的最小代数闭域. 这一构造最熟悉的应用是 $\overline{\mathbb{R}}=\mathbb{C}$. 把此方法应用到任意域, 就能得到许多新的代数闭域例子, 尽管它们并不特别熟悉, 如 $\overline{\mathbb{Q}}$, $\overline{\mathbb{F}}_{p}$ 等. 不熟悉这些内容的读者, 每当遇到 $K$ 是代数闭域的假设时, 不妨想一想较熟悉的情形 $K=\mathbb{C}$.

在讨论了何谓代数闭域以后, 现在可以叙述零点定理了. 它的证明留到第~\ref{chap:nullstellensatz-proof}~章.

\begin{Theorem}{零点定理}\leavevmode\label{thm:nullstellensatz}
  设 $K$ 是代数闭域. 则对任意理想 $I\subseteq K[x_{1},\ldots,x_{n}]$, 有
  \[
    \mathcal{I}(\mathcal{V}(I))=\sqrt{I}.
  \]
\end{Theorem}

包含关系 $\sqrt{I}\subseteq\mathcal{I}(\mathcal{V}(I))$ 在任意域上都成立, 可以直接由定义证明 (Exercise~\ref{ex:radical-contained-in-vanishing-ideal}). 另一方向的包含关系需要相当多的工作, 而且 $K$ 为代数闭域这一假设是必要的.

零点定理帮助我们回答第~\ref{sec:affine-varieties}~节末尾提出的问题: 在怎样的定义域和陪域上, $\mathcal{V}$ 算子是双射? 在代数闭域上, 答案是 $\mathcal{V}$ 给出根理想与仿射簇之间的一个双射. 这是代数与几何之间的精确词典中第一个关键例子.

\begin{Corollary}{根理想与仿射簇}\leavevmode\label{cor:radical-ideals-affine-varieties-bijection}
  若 $K$ 是代数闭域, 则
  \[
    \mathcal{V}:\{K[x_{1},\ldots,x_{n}]\text{ 的根理想}\}
      \longrightarrow\{\mathbb{A}^{n}\text{ 中的仿射簇}\}
  \]
  是一个双射, 其逆为 $\mathcal{I}$.
\end{Corollary}

\begin{Proof}
  由于一个函数是双射当且仅当它有逆函数, 只需证明 $\mathcal{I}$ 是 $\mathcal{V}$ 的逆. 换句话说, 我们必须证明对任意仿射簇 $X\subseteq\mathbb{A}^{n}$, 有
  \begin{equation}\label{eq:vanishing-of-affine-variety-ideal}
    \mathcal{V}(\mathcal{I}(X))=X
  \end{equation}
  并且对任意根理想 $I\subseteq K[x_{1},\ldots,x_{n}]$, 有
  \begin{equation}\label{eq:ideal-of-radical-vanishing-set}
    \mathcal{I}(\mathcal{V}(I))=I.
  \end{equation}
  等式 \eqref{eq:vanishing-of-affine-variety-ideal} 是 Proposition~\ref{prop:compose-vanishing-operators} 的第一部分, 因此无须对 $K$ 作任何假设. 为证 \eqref{eq:ideal-of-radical-vanishing-set}, 注意
  \[
    \mathcal{I}(\mathcal{V}(I))=\sqrt{I}=I,
  \]
  其中第一个等号来自零点定理, 第二个等号来自根理想的刻画, 即根理想恰好是等于自身的根的理想 (Corollary~\ref{cor:characterize-radical-ideals}).
\end{Proof}

零点定理不仅给出仿射簇和根理想之间的双射, 也是计算消没理想的有用工具. 特别地, 判断一个理想是否为根理想, 常常比判断它是否为消没理想容易得多. 以下几个例子说明这一点.

\begin{Example}{$\mathcal{I}(\mathcal{V}(y-x^{2}))=\langle y-x^{2}\rangle$}\leavevmode\label{ex:parabola-nullstellensatz}
  设 $K$ 是代数闭域, 并考虑仿射簇
  \[
    X=\mathcal{V}(y-x^{2})\subseteq\mathbb{A}^{2}.
  \]
  由于每个代数闭域都是无限域, 由 Example~\ref{ex:vanishing-ideal-of-parabola} 已知 $\mathcal{I}(X)=\langle y-x^{2}\rangle$. 不过, 该例中的论证颇费周折, 且只针对特定的多项式 $y-x^{2}$. 对代数闭域, 有一个适用于所有不可约多项式的简单得多的论证.

  具体来说, 已知 $y-x^{2}\in K[x,y]$ 是不可约多项式, 可知 $\langle y-x^{2}\rangle\subseteq K[x,y]$ 是素理想 (Proposition~\ref{prop:irreducible-polynomial-characterization}), 因而是根理想 (Proposition~\ref{prop:prime-ideals-are-radical}). 于是由零点定理,
  \[
    \mathcal{I}(X)=\mathcal{I}(\mathcal{V}(y-x^{2}))
      =\sqrt{\langle y-x^{2}\rangle}=\langle y-x^{2}\rangle.
  \]
\end{Example}

\begin{Example}{$f$ 无平方因子时 $\mathcal{I}(\mathcal{V}(f))=\langle f\rangle$}\leavevmode\label{ex:square-free-polynomial-nullstellensatz}
  推广上一例可知, 只要 $K$ 是代数闭域, 且 $f\in K[x_{1},\ldots,x_{n}]$ 不被任何非零多项式的平方整除, 就有
  \[
    \mathcal{I}(\mathcal{V}(f))=\sqrt{\langle f\rangle}=\langle f\rangle.
  \]
  第一个等号来自零点定理, 第二个等号来自 Proposition~\ref{prop:radicals-of-principal-ideals}. 特别地, 若 $f$ 不可约, 则 $\mathcal{V}(f)$ 的消没理想就是 $\langle f\rangle$.

  注意, 当 $K$ 不是代数闭域时, 本例的结论不成立. 例如, 考虑不可约多项式 $x^{2}+1\in\mathbb{R}[x]$. 由于它没有零点,
  \[
    \mathcal{I}(\mathcal{V}(x^{2}+1))=\mathcal{I}(\varnothing)
      =K[x]\supsetneq\langle x^{2}+1\rangle.
  \]
\end{Example}

对我们将继续在代数与几何之间建立的丰富词典而言, 根理想与仿射簇之间的双射仅仅触及表面. 借用语言来类比, Corollary~\ref{cor:radical-ideals-affine-varieties-bijection} 中的双射应当看作两种语言之间对名词的翻译. 这种翻译或许能让我们进行十分简单的对话, 但若要充分利用语言的丰富性, 还应翻译动词, 形容词, 副词等等.

在接下来的三章中, 我们会继续建立代数与几何之间的词典, 并始终假设零点定理成立. 在这个过程中, 我们将有机会引入一些有用的新代数概念. 到第~\ref{chap:nullstellensatz-proof}~章, 当我们有了更扎实的代数基础, 对代数与几何之间的词典有了更全面的认识之后, 再回头证明零点定理.

\begin{BookGroundFieldNote}
  \noindent\textbf{\color{red!45!black}关于基域 $K$ 的假设}\par\smallskip
  零点定理是代数几何的主干. 因此, \textbf{在本书余下部分, 若无特别说明, 均假设 $K$ 是代数闭域.} 我们所建立的许多定义和结果在一般的域上依然成立. 但另一些需要稍作调整, 还有一些在基域不是代数闭域时完全不成立. 为了帮助读者理解我们的假设, 对代数闭性不可缺少的结果, 我们常用 $K=\mathbb{R}$ 的情形来说明反例.

  虽然代数几何的核心结果在实数域上不成立, 因为实数域不是代数闭域, 但我们对仿射簇的很多几何直觉来自观察多项式在 $\mathbb{R}$ 上的解. 的确, 第~\ref{sec:vanishing-set-operator}~节中每一幅关于零点集的几何图像都画的是 $\mathbb{R}$ 上的仿射簇. 作为代数几何研究者, 我们需要学会把在 $\mathbb{R}$ 上的知识和直觉用作启发, 同时不被那一特殊情形中因 $\mathbb{R}$ 并非代数闭域而出现的现象所误导.

  接下来, 虽然总是假设基域是代数闭域, 但我们仍会通过观察 $\mathbb{R}$ 上的解来讨论和描绘簇的例子, 并继续使用我们多年来研究这些集合时熟悉的词语. 例如, 我们把簇 $\mathcal{V}(y-x^{2})\subseteq\mathbb{A}^{2}$ 称为\emph{抛物线}, 把 $\mathcal{V}(x^{2}+y^{2}+z^{2}-1)\subseteq\mathbb{A}^{3}$ 称为\emph{单位球面}, 尽管在一般的域上, 这些簇可能与\emph{抛物线}和\emph{球面}这些词在我们脑海中唤起的几何图形并不十分相似. 因为 $\mathbb{R}$ 是代数闭域 $\mathbb{C}$ 的一个子集, 读者不妨始终假设 $K=\mathbb{C}$. 在这种情形下, 例子中所画的 $\mathbb{R}$ 上的图像是 $\mathbb{C}$ 上完整解集的一个子集. 这些图像并不能给出全貌, 但至少可让我们由一瞥, 或者说一个切片, 略窥这个簇的性质.

  域的另一个重要属性是它的\textbf{特征}. 回顾一下, $K$ 的特征记为 $\operatorname{char}(K)$, 是使得下式成立的最小正整数 $p$:
  \[
    \underbrace{1+\cdots+1}_{p}=0\in K.
  \]
  如果不存在这样的 $p$, 如 $\mathbb{Q}$, $\mathbb{R}$ 和 $\mathbb{C}$ 的情形, 就说这个域的特征为 $0$. 本书中所有结果都对一般的代数闭域成立. 不过, 在举例时, 我们可能希望避开有限多个特征, 因为在那些特征下, 某些类型的多项式可能表现异常. 例如, 当 $\operatorname{char}(K)\ne2$ 时, 多项式 $x^{2}+y^{2}-1$ 不可约, 但如果 $1+1=0\in K$, 那么
  \[
    x^{2}+y^{2}-1=(x+y+1)^{2}.
  \]
  为了简明起见, 在具体例子中我们常假设 $K=\mathbb{C}$, 而不列出例外的特征, 尽管例子所给出的结论通常能推广到一般的代数闭域, 只有有限多个特征例外.
\end{BookGroundFieldNote}

\BookExerciseHeading
\begin{BookExercises}
  \item\label{ex:algebraically-closed-field-infinite} 证明任意代数闭域都是无限域. (提示: 如果 $K=\{a_{1},\ldots,a_{r}\}$ 是有限集, 你能在 $K[x]$ 中构造一个在 $K$ 中没有零点的多项式吗?)
  \item\label{ex:factorization-over-algebraically-closed-field} 设 $K$ 是代数闭域, 并设 $f\in K[x]$ 是次数为 $d$ 的多项式. 证明存在 $a_{0},a_{1},\ldots,a_{d}\in K$, 不要求彼此不同, 使得
    \[
      f=a_{0}(x-a_{1})\cdots(x-a_{d}).
    \]
  \item\label{ex:radical-contained-in-vanishing-ideal} 设 $K$ 是任意域, $I\subseteq K[x_{1},\ldots,x_{n}]$ 是一个理想. 证明零点定理的一个包含关系不需要假设 $K$ 是代数闭域:
    \[
      \sqrt{I}\subseteq\mathcal{I}(\mathcal{V}(I)).
    \]
  \item\label{ex:nullstellensatz-requires-algebraically-closed-field} 证明零点定理在任何不是代数闭域的域上均不成立.
  \item\label{ex:weak-nullstellensatz} 设 $K$ 是代数闭域, $I\subseteq K[x_{1},\ldots,x_{n}]$ 是一个理想. 假设零点定理成立, 证明 $\mathcal{V}(I)=\varnothing$ 当且仅当 $I=K[x_{1},\ldots,x_{n}]$. (这个结果通常称为弱零点定理.)
  \item\label{ex:compute-ideals-by-nullstellensatz} 假设零点定理成立, 求以下各簇的 $\mathcal{I}(X)$.
    \begin{enumerate}[label=(\alph*)]
      \item $X=\mathcal{V}(x^{2}-y^{3},x^{2}+y^{3})\subseteq\mathbb{A}^{2}_{\mathbb{C}}$
      \item $X=\mathcal{V}(x)\cup\mathcal{V}(y-z)\subseteq\mathbb{A}^{3}_{\mathbb{C}}$
      \item $X=\mathcal{V}((x^{2}y^{2}+yz^{2}+x^{3}z^{2})(x^{2}+y^{2}+1))\subseteq\mathbb{A}^{3}_{\mathbb{C}}$
    \end{enumerate}
  \item\label{ex:parametrized-radical-ideal-parabola-intersection} 设 $K$ 是代数闭域. 对每个 $a\in K$, 定义
    \[
      X_{a}=\mathcal{V}(y-x^{2},y-a)\subseteq\mathbb{A}^{2}_{K}.
    \]
    \begin{enumerate}[label=(\alph*),ref=(\alph*)]
      \item\label{item:parabola-intersection-radical-values} $a$ 取哪些值时有 $\mathcal{I}(X_{a})=\langle y-x^{2},y-a\rangle$? 等式不成立时, $\mathcal{I}(X_{a})$ 是什么?
      \item 对 $a\in\mathbb{R}$ 取几个有代表性的值, 在 $\mathbb{R}$ 上画出仿射簇 $X_{a}$. 你能从几何上解释 \ref{item:parabola-intersection-radical-values} 中等式成立的 $a$ 值与等式不成立的 $a$ 值之间的差别吗?
    \end{enumerate}
  \item\label{ex:unit-circle-characteristic-two} 假设 $\operatorname{char}(K)=2$. 证明 $x^{2}+y^{2}-1=(x+y+1)^{2}\in K[x]$.
  \item\label{ex:frobenius-additivity-for-polynomials} 推广上一例, 假设对某个素数 $p\in\mathbb{Z}$, 有 $\operatorname{char}(K)=p$. 证明对任意 $f_{1},\ldots,f_{m}\in K[x_{1},\ldots,x_{n}]$, 都有
    \[
      \left(\sum_{i=1}^{m}f_{i}\right)^{p}
        =\sum_{i=1}^{m}f_{i}^{p}\in K[x_{1},\ldots,x_{n}].
    \]
\end{BookExercises}
